\section{Methodology}
\label{sec:methodology}
\subsection{Study sites}

Experiments are conducted at two meteorological stations in Colombia that represent
climatologically contrasting regimes (Table~\ref{tab:sites}).

\begin{table}[t]
  \centering
  \caption{Study site metadata.}
  \label{tab:sites}
  \begin{tabular}{lcc}
    \toprule
    & \textbf{El Paso (César)} & \textbf{Uniandes (Bogotá)} \\
    \midrule
    Latitude  & $9.737^{\circ}\text{N}$ & $4.602^{\circ}\text{N}$ \\
    Longitude & $73.695^{\circ}\text{W}$ & $74.066^{\circ}\text{W}$ \\
    Elevation & 50~m a.s.l.      & 2\,600~m a.s.l.  \\
    Climate   & Tropical semi-arid & Equatorial mountain \\
    Avg.\ GHI & ${\sim}$5.5~kWh/m²/day & ${\sim}$4.2~kWh/m²/day \\
    \bottomrule
  \end{tabular}
\end{table}

\textbf{El Paso} is located in the lowland Caribbean region of C{\'e}sar department,
characterised by high annual irradiance, predominantly clear-sky conditions, and
dry-season cloud cover driven by the Inter-Tropical Convergence Zone (ITCZ).

\textbf{Universidad de los Andes} (referred to as Uniandes) is an instrumented
rooftop station in the Colombian Andean city of Bogot{\'a}. At 2\,600~m above sea
level, it experiences the equatorial bimodal rainfall regime with intense convective
cloud formation typically peaking in the early afternoon, producing sharp intraday
GHI drops.

\subsection{Data}

\subsubsection{Ground measurements}

GHI is measured at 10-minute resolution and stored in time-zone-aware (UTC) Parquet
files. Missing values account for less than 2\% of the record at both sites and are
excluded during dataset construction.

\subsubsection{Satellite data and patch extraction}
\label{sssec:satellite}

We use the GOES-16 ABI Level-2 Multi-Band Cloud and Moisture Imagery Product
(ABI-L2-MCMIPF), which provides all 16 ABI spectral channels at a unified
10-minute cadence and ${\sim}2$~km spatial resolution in the visible bands.
Raw hourly ABI-L2-MCMIPF granules are pre-processed into compact per-hour
archives: each archive is a compressed NumPy array of shape
$(6, 16, P, P)$, holding the six 10-minute slots within that hour, the
$C = 16$ spectral channels, and a $P \times P = 16 \times 16$ pixel patch
($P = 16$) centred on the target site, in units consistent with each band's
native reflectance or brightness temperature. At inference/training time, a
single timestamp is resolved to its $(C, P, P) = (16, 16, 16)$ frame by
selecting the corresponding hourly archive and 10-minute slot, so no
redundant decoding of the original satellite granules is required.

\begin{figure}[t]
  \centering
  \includegraphics[width=\columnwidth]{fig7_patch_snapshot}
  \caption{Example $16\times16$ satellite patch (a single ABI reflectance
           channel, one 10-minute slot) at each site, same daytime hour.
           El~Paso shows a spatially coherent cloud edge typical of the
           site's predominantly clear-sky regime; Uniandes shows more
           fragmented, patchy structure consistent with locally-forming
           convective cloud cover (Section~\ref{sec:discussion}).}
  \label{fig:patch-snapshot}
\end{figure}

Figure~\ref{fig:patch-snapshot} shows an example patch at each site for
orientation.

\subsubsection{Manifests, sequence construction, and temporal split}
\label{sssec:manifest}

For the deep-learning models (ResNet-LSTM, GraphSAGE-LSTM, FlatMLP), forecast
samples are constructed as pairs $(X, y)$ where the input
$X \in \R^{L \times C \times P \times P}$ is a sequence of $L = 24$ satellite
patch frames covering the most recent 4 hours of observation (10-minute
cadence), and the target $y \in \R$ is the ground-measured GHI at the site
$H$ steps ahead ($H \in \{6, 18, 36\}$ steps, corresponding to 1, 3, and
6 hours).

A pre-computed \textbf{manifest} (one Parquet table per site and horizon)
enumerates every valid forecast sample: for each candidate target timestamp it
stores the label $y$, the list of $L = 24$ history timestamps
\texttt{history\_ts}, and a coverage flag confirming that (i) all $L$ history
steps exist in the ground-truth series and (ii) a satellite archive is
available for each of them. Samples for which any of the $L$ history frames
is missing on disk are filtered out at load time, so every batch the models
see contains a complete 4-hour satellite history. Two thin
\texttt{Dataset} wrappers consume the same manifest and patch archives but
reshape the retrieved frames differently for the two model families:

\begin{itemize}
  \item \textbf{\texttt{PatchSeqDataset}} (used by ResNet-LSTM and FlatMLP)
        returns each sample as an image-like sequence
        $X \in \R^{L \times C \times P \times P}$ --- i.e., $L$ stacked
        $(16, 16, 16)$ tensors preserving the 2-D pixel grid, ready to be
        consumed by 2-D convolutions.
  \item \textbf{\texttt{GraphSeqDataset}} (used by GraphSAGE-LSTM) returns
        each sample as a node-feature sequence $X \in \R^{L \times N \times C}$
        with $N = P^2 = 256$: every frame is transposed and reshaped from
        $(C, P, P)$ to $(P^2, C)$, so that each of the 256 pixels becomes a
        graph node carrying a $C = 16$-dimensional feature vector (its
        16-channel spectral signature at that location).
\end{itemize}

To prevent data leakage, train, validation, and test sets are defined by
non-overlapping calendar blocks (Table~\ref{tab:splits}), and the manifest
coverage check, splitting boundaries, and patch archives are shared across
all three deep-learning models, ensuring that ResNet-LSTM, GraphSAGE-LSTM,
and FlatMLP are trained and evaluated on \emph{exactly} the same samples.

\begin{table}[t]
  \centering
  \caption{Temporal data splits.}
  \label{tab:splits}
  \setlength{\tabcolsep}{4pt}
  \begin{tabular}{llccc}
    \toprule
    \textbf{Site} & \textbf{Split} & \textbf{Start} & \textbf{End} & $n$ (approx.) \\
    \midrule
    \multirow{3}{*}{El Paso}
      & Train & 2022-03-01 & 2023-06-30 & 69\,984 \\
      & Val   & 2023-07-01 & 2023-10-31 & 17\,568 \\
      & Test  & 2023-11-01 & 2024-03-07 & 18\,288 \\
    \midrule
    \multirow{3}{*}{Uniandes}
      & Train & 2023-09-01 & 2024-09-30 & 56\,850 \\
      & Val   & 2024-10-01 & 2024-12-31 & 13\,104 \\
      & Test  & 2025-01-01 & 2025-03-28 & 12\,384 \\
    \bottomrule
  \end{tabular}
\end{table}

\subsection{Models}

\subsubsection{Persistence baseline}

The persistence model assumes irradiance is unchanged over the forecast horizon:
$\hat{y}(t+H) = \text{GHI}(t)$. It is optimal under no-change dynamics and
serves as the standard benchmark for short-term solar forecasting.

\subsubsection{SARIMA: classical statistical baseline}
\label{ssec:sarima}

As a non-deep-learning reference, we fit a seasonal ARIMA model,
SARIMA$(2,1,2)(1,1,1)_{24}$ --- implemented via \texttt{statsmodels}'
\texttt{SARIMAX} class with no exogenous regressors, i.e.\ a plain SARIMA
model rather than a SARIMAX one in the methodological sense --- directly on
the raw hourly GHI series at each site. A single order was selected via AIC/BIC grid search (documented
separately) and applied uniformly at both sites, rather than re-selected
per site, to keep the classical-statistical reference comparable across
sites; no clear-sky normalisation is applied, so the model targets the same
physical quantity, GHI in W/m$^2$, as every other model.

Evaluation is \emph{rolling-origin}: model parameters are estimated once on
the training period and then held fixed, and the underlying state-space
(Kalman) filter is advanced observation-by-observation through the
validation and test periods with no re-estimation. At each test forecast
origin $t$, a genuine $H$-hour-ahead forecast of $t+H$ is issued from the
filtered state that incorporates all observations up to and including $t$
--- exactly the information persistence and the deep-learning models are
allowed to condition on, and never the future. Because a seasonal period of $m=144$
(10-minute native resolution) is computationally intractable over the
multi-year training window, the satellite-resolution series is first
resampled to hourly cadence ($m=24$; missing hours are handled natively by
the filter as missing observations), and SARIMA is therefore evaluated only
on the subset of test timestamps that fall on an exact hour, with the
persistence reference \emph{recomputed on that same coarser subset} so that
SARIMA's own skill score is internally consistent. \textbf{This makes
SARIMA's $\rmseday$ and $\skillday$ not directly comparable to the
full 10-minute-resolution metrics reported for persistence and the
deep-learning models elsewhere in this paper} --- both because of the
coarser evaluation subset and because SARIMA is fit on, and therefore
forecasts, hourly-\emph{mean} GHI while being scored against the same
instantaneous 10-minute targets as every other model; it is included in
Figs.~\ref{fig:skill}--\ref{fig:rmse} as a classical-statistical reference
point, visually flagged (hatched bars), rather than in Table~\ref{tab:main}.

\subsubsection{ResNet-LSTM: convolutional spatial encoder}
\label{ssec:resnet}

ResNet-LSTM is an encoder--decoder that treats each satellite frame as a
conventional multi-channel image and lets a convolutional network learn its
spatial representation, before a recurrent network reasons over the sequence
of representations (Figure~\ref{fig:arch-resnet}).

\begin{figure*}[t]
  \centering
  \includegraphics[width=0.85\textwidth]{fig_arch_resnet_lstm}
  \caption{ResNet-LSTM architecture. Each of the $L=24$ satellite frames is
           encoded independently by a shared lightweight ResNet, producing a
           per-frame embedding $\mathbf{e}_\ell$; the resulting sequence is
           processed by an LSTM and a final MLP head to produce the point
           forecast $\hat{y}$.}
  \label{fig:arch-resnet}
\end{figure*}

\textbf{Spatial encoder (per frame).} Each frame $X_\ell \in \R^{C \times P
\times P}$ ($\ell = 1,\ldots,L$) is passed independently through a lightweight
ResNet \cite{he2016}: a $3\times3$ convolutional stem, followed by three
residual stages of \texttt{BasicBlock} units (two $3\times3$ convolutions with
batch normalisation and a skip connection). The first stage preserves the
$16\times16$ spatial resolution, while the second and third stages apply
stride-2 downsampling, taking the spatial grid through
$16 \to 8 \to 4$. A global-average-pooling layer then collapses the
$4\times4$ feature map to a single vector, which a final linear layer projects
to a $d_e$-dimensional frame embedding $\mathbf{e}_\ell \in \R^{d_e}$
(\texttt{emb\_dim}, tuned in $\{64,128,192,256\}$). The encoder's channel
width is controlled by \texttt{base}~$\in\{16,24,32,48\}$ (the three stages
use \texttt{base}, $2\times$\texttt{base}, and $4\times$\texttt{base}
channels respectively); the resulting parameter count ranges from
approximately 50k to 500k depending on the configuration. Because the same
encoder weights are applied to every one of the $L = 24$ frames, the model is
translation-equivariant and shares a single spatial template across the whole
4-hour history.

\textbf{Temporal encoder.} The $L$ frame embeddings are normalised
(\texttt{LayerNorm}) and stacked into a sequence
$(\mathbf{e}_1,\ldots,\mathbf{e}_L) \in \R^{L \times d_e}$, which a
single- or two-layer LSTM \cite{hochreiter1997} (\texttt{hidden\_t}, \texttt{n\_lstm\_layers})
processes step by step; only its final hidden state $\mathbf{h}_L \in
\R^{d_t}$ is retained. A two-layer MLP head
(\texttt{Linear} $\to$ \texttt{ReLU} $\to$ \texttt{Dropout} $\to$
\texttt{Linear}) maps $\mathbf{h}_L$ to the scalar point forecast $\hat{y}$.

\subsubsection{GraphSAGE-LSTM: graph-based spatial encoder}
\label{ssec:gsage}

GraphSAGE-LSTM shares the same encoder--decoder skeleton as ResNet-LSTM ---
spatial encoder per frame, then LSTM over the embedding sequence, then MLP
head --- but \emph{replaces the convolutional encoder with a graph neural
network} that treats individual pixels, rather than the whole frame, as the
basic spatial unit (Figure~\ref{fig:arch-gsage}).

\begin{figure*}[t]
  \centering
  \includegraphics[width=0.85\textwidth]{fig_arch_graphsage_lstm}
  \caption{GraphSAGE-LSTM architecture. Each satellite frame is reshaped into
           $N=256$ pixel-nodes and processed by $K$ stacked GraphSAGE layers
           over a $k$-NN graph (fixed unweighted baseline or tuned
           distance-weighted variant, Section~\ref{ssec:gsage}); a mean
           readout produces the per-frame graph embedding $\mathbf{g}_\ell$,
           which feeds the same LSTM and MLP head as ResNet-LSTM.}
  \label{fig:arch-gsage}
\end{figure*}

\textbf{From pixels to a graph.} Each $(C, P, P)$ frame is reshaped (by
\texttt{GraphSeqDataset}, Section~\ref{sssec:manifest}) into $N = P^2 = 256$
nodes, one per pixel, each carrying its native $C = 16$-channel spectral
vector as a node feature --- no convolution or pooling is applied beforehand.
A graph $G = (V, E, \mathbf{w})$ is built once over the $16\times16$ pixel
grid and reused for every frame and every sample; the edge index and weights
are precomputed once, registered as model buffers, and replicated across the
batch dimension, so the same mini-batch \texttt{DataLoader} and training loop
used for ResNet-LSTM apply unchanged. We compare two graph-construction
variants, isolating the contribution of \emph{tuned spatial connectivity}
to forecasting skill (ablation reported in Section~\ref{sec:results}):

\begin{itemize}
  \item \textbf{Fixed unweighted graph (baseline).} Each node connects to its
        8 immediate grid neighbours (cardinal and diagonal) with unit edge
        weight, the standard local-connectivity topology used in early graph
        convolutional networks applied to image-like data.
  \item \textbf{Tuned weighted $k$-NN graph (proposed).} Each node connects
        to its $k$ closest pixels by Euclidean distance in 2-D pixel
        coordinates ($k \in \{4,8,12,16\}$, tuned per site/horizon,
        Section~\ref{ssec:optuna}), and every edge is assigned a fixed
        inverse-distance weight $w_{vu} = 1/d(v,u)$
        (Eq.~\ref{eq:edge_weight} below) --- the edge weights themselves are
        a deterministic function of pixel geometry, not a trainable
        parameter; only the neighbourhood size $k$ is selected by Optuna.
        Larger $k$ values progressively add diagonal, two-step, and
        knight-move neighbours, so $k$ effectively controls the spatial
        receptive field of a single graph layer. Unlike the fixed baseline
        or a convolutional filter, this topology lets the \emph{trainable}
        GraphSAGE weights $\mathbf{W}^{(k)}$ (Eq.~\ref{eq:sage}) apply
        \emph{different} effective aggregation to differently-located pixels
        (e.g.\ patch centre vs.\ periphery), rather than sliding one shared
        template across the image, and lets Optuna select the neighbourhood
        size that best matches each site's cloud spatial scale.
\end{itemize}

Both variants are trained and reported for Uniandes. For El~Paso, a
satellite-timestamp alignment bug (Section~\ref{sec:discussion},
Limitations) invalidated all pre-fix results, and by the time this was
found, the fixed-graph training code path had already been superseded by
the tuned-graph one described above; rather than reinstate discontinued
code to reproduce a fixed-graph result for El~Paso, we report El~Paso's
GraphSAGE-LSTM using only the tuned-graph variant once its corrected
retraining completes, with no fixed-vs-tuned ablation at that site.

\begin{equation}
  w_{vu} = \frac{1}{d(v,u)}, \quad (v, u) \in E.
  \label{eq:edge_weight}
\end{equation}

\textbf{Spatial encoder (per frame).} GraphSAGE \cite{hamilton2017} updates
every node's representation by aggregating its neighbourhood. With edge
weights, the neighbourhood mean becomes a distance-weighted average, and the
update concatenates each node's own (self) representation with that
weighted-neighbour summary before a shared linear projection and ReLU:
\begin{equation}
  \mathbf{h}_v^{(k)} = \sigma\!\left(
    \mathbf{W}^{(k)} \left[ \mathbf{h}_v^{(k-1)} \Vert
    \frac{\sum_{u \in \mathcal{N}(v)} w_{vu}\,\mathbf{h}_u^{(k-1)}}
         {\sum_{u \in \mathcal{N}(v)} w_{vu}} \right]
  \right)
  \label{eq:sage}
\end{equation}
where $\mathcal{N}(v)$ denotes the neighbours of $v$ under whichever graph
variant is in use, $\Vert$ is concatenation, and $\sigma$ is ReLU. Two further
implementation choices are left to Optuna rather than fixed a priori
(\texttt{input\_bn}, \texttt{concat\_agg} in Table~\ref{tab:hpspace}): whether
node features are passed through a BatchNorm1d layer before the first SAGE
layer, and whether the aggregation step uses the concatenation of
Eq.~\ref{eq:sage} or an additive alternative,
$\mathbf{h}_v^{(k)} = \sigma(\mathbf{W}^{(k)}_\text{self}\mathbf{h}_v^{(k-1)} +
\mathbf{W}^{(k)}_\text{nei}\,\overline{\mathbf{h}}_u^{(k-1)})$, which halves
the number of parameters per layer relative to concatenation. For the
fixed unweighted baseline, $w_{vu} = 1$ for all $(v,u) \in E$, reducing
Eq.~\ref{eq:sage} to the standard unweighted SAGE neighbourhood mean; the two
graph variants therefore share the same aggregation layer and differ only in
$E$ and $\mathbf{w}$, isolating the effect of graph construction from the
rest of the architecture. The same layer (with shared weights
$\mathbf{W}^{(k)}$, output width \texttt{hidden\_g}) is applied to all 256
nodes simultaneously and stacked $K \in \{2,3,4\}$ times
(\texttt{n\_sage\_layers}); each additional layer extends a node's effective
receptive field by one more $k$-NN hop. A final \emph{mean readout} over all
$N = 256$ nodes collapses the per-pixel representations into a single
graph-level (frame) embedding $\mathbf{g}_\ell \in \R^{d_g}$, with
$d_g = $~\texttt{hidden\_g}~$\in \{64,96,128,192,256\}$ --- the GraphSAGE
analogue of the ResNet's $d_e$-dimensional frame embedding.

\textbf{Temporal encoder.} As in ResNet-LSTM, the $L$ graph-level embeddings
are normalised and stacked into $(\mathbf{g}_1,\ldots,\mathbf{g}_L) \in
\R^{L \times d_g}$ and processed by a single- or two-layer LSTM whose final
hidden state feeds an identical MLP head producing $\hat{y}$. The two
architectures are thus directly comparable: they differ \emph{only} in how
the per-frame spatial embedding is produced --- convolution-and-downsampling
over a regular grid (ResNet) versus weighted message-passing over a tuned
$k$-NN pixel graph at full $256$-node resolution, with no spatial downsampling
(GraphSAGE) --- while sharing the same temporal model, output head, training
protocol, and data pipeline.

\subsubsection{FlatMLP: spatial-ablation baseline}
\label{ssec:mlp}

To isolate the contribution of \emph{spatial} encoding from that of temporal
modelling, FlatMLP discards all pixel-level spatial structure
(Figure~\ref{fig:arch-mlp}). For every
frame $X_\ell \in \R^{C \times P \times P}$, a global spatial average pool
collapses the $P \times P$ grid to its per-channel mean, $\bar{\mathbf{x}}_\ell
\in \R^{C}$ --- i.e.\ each frame is reduced to a 16-dimensional ``patch-mean
spectral signature'' with no notion of pixel location. The $L = 24$ resulting
vectors are concatenated into a single flat input
$\bar{\mathbf{x}} \in \R^{L \times C} = \R^{384}$ ($24 \times 16$). This
vector is passed through a fully-connected feed-forward network of
\texttt{n\_layers} blocks, each consisting of \texttt{Linear} $\to$
\texttt{LayerNorm} $\to$ \texttt{ReLU} $\to$ \texttt{Dropout}
(\texttt{hidden\_dim} units), followed by a final \texttt{Linear} layer that
outputs the scalar $\hat{y}$ directly --- every unit in one layer is connected
to every unit in the next (a standard dense MLP), with no recurrence, no
convolution, and no graph structure. Because FlatMLP receives exactly the
same $L\times C \times P \times P$ input tensors (via \texttt{PatchSeqDataset})
and is trained with the same protocol, its skill gap relative to ResNet-LSTM
and GraphSAGE-LSTM provides an approximate estimate of how much of their
forecasting skill stems from \emph{spatial} encoding of the satellite patch,
as opposed to simple per-frame statistics processed over time. Note that
FlatMLP also replaces the shared LSTM with a feed-forward network over the
concatenated sequence, so its skill gap reflects the joint loss of spatial
structure \emph{and} of the recurrent temporal model, not spatial encoding
in isolation (Section~\ref{sec:discussion}, Limitations).

\begin{figure}[t]
  \centering
  \includegraphics[width=\columnwidth]{fig_arch_flatmlp}
  \caption{FlatMLP architecture. Spatial encoding is reduced to a per-frame
           global average pool (no convolution, graph, or recurrence over
           space), so any skill gap relative to ResNet-LSTM or GraphSAGE-LSTM
           is directly attributable to spatial encoding of the satellite
           patch.}
  \label{fig:arch-mlp}
\end{figure}

\subsection{Hyperparameter optimisation}
\label{ssec:optuna}

For the three deep-learning models (ResNet-LSTM, GraphSAGE-LSTM, and FlatMLP)
we use Optuna \cite{akiba2019} with the TPE sampler and Median pruner.
The fixed-graph GraphSAGE-LSTM baseline and the corresponding ResNet-LSTM
configuration use a budget of 20 trials per configuration (site $\times$
horizon $\times$ seed) across four random seeds
$\{42,1,7,13\}$;\footnote{One fixed-graph configuration (GraphSAGE-LSTM,
Uniandes, 6-hour horizon, seed~1) was run with 50 trials during an early
budget experiment; its result is included in the four-seed statistics
as obtained.} FlatMLP
uses a larger budget of 50 trials given its lower per-trial cost. The
extended search space used for the weighted $k$-NN ablation
(Table~\ref{tab:hpspace}) increases per-trial cost substantially --- larger
embeddings, optionally deeper LSTMs, and a tunable graph rebuilt per trial
--- so this variant instead uses a larger per-seed budget of 75--100 trials
(100 in the initial Uniandes configurations; 75 under the standardised
protocol used for all subsequent runs)
across a reduced set of two random seeds $\{42,1\}$, trading seed count for
per-seed search coverage of the most consequential hyperparameters (graph
connectivity, model capacity); results are reported with this caveat
wherever the two protocols are compared (Section~\ref{sec:results}).
El~Paso's corrected retraining (Section~\ref{sec:results} note;
Section~\ref{sec:discussion}, Limitations) follows this extended-space
two-seed protocol for \emph{both} architectures, not only the GraphSAGE-LSTM
ablation, and the affected tables state this alongside the numbers.

Table~\ref{tab:hpspace} lists the \emph{extended} search space. Relative to
the initial 20-trial spaces behind the fixed-graph four-seed results, it
adds a categorical L1 regularisation coefficient (\texttt{l1\_reg}) for both
architectures, to penalise large weights in high-capacity configurations,
and --- for ResNet-LSTM --- stacked recurrent processing via
\texttt{n\_lstm\_layers}~$\in\{1,2\}$ (GraphSAGE-LSTM's initial space
already tuned its LSTM depth).
For \textbf{ResNet-LSTM} specifically, the embedding dimension
(\texttt{emb\_dim}) and temporal hidden size (\texttt{hidden\_t}) are extended
to 256 to increase capacity.
For \textbf{GraphSAGE-LSTM}, the number of nearest neighbours
($k \in \{4, 8, 12, 16\}$) used to build the weighted pixel graph is
additionally tunable, allowing the study to determine the optimal spatial
connectivity for each site-horizon combination, together with the
\texttt{input\_bn} and \texttt{concat\_agg} choices described in
Section~\ref{ssec:gsage}. The fixed-graph baseline
(Section~\ref{ssec:gsage}) is tuned over the initial space (the same
architectural choices minus \texttt{k\_neighbors}, whose topology is fixed
a priori, and minus \texttt{l1\_reg}), and serves as the
ablation reference for the value of tuned spatial connectivity.
For \textbf{FlatMLP}, the study searches over MLP depth, hidden dimension,
dropout, learning rate, and weight decay (50 trials).

\begin{table*}[t]
  \centering
  \caption{Optuna hyperparameter search spaces (extended protocol). The
           \texttt{l1\_reg} coefficient is new relative to the initial
           20-trial spaces for both architectures; \texttt{n\_lstm\_layers}
           is new for ResNet-LSTM only. The \texttt{k\_neighbors}
           parameter determines the weighted $k$-NN graph built per trial for
           GraphSAGE-LSTM. \texttt{input\_bn} toggles the input BatchNorm1d
           described below; \texttt{concat\_agg} toggles between the
           concatenation-based aggregation of Eq.~\ref{eq:sage} and an additive
           alternative (Section~\ref{ssec:gsage}) --- both are tuned per trial
           rather than fixed a priori.}
  \label{tab:hpspace}
  \setlength{\tabcolsep}{5pt}
  \begin{tabular}{lll}
    \toprule
    \textbf{Hyperparameter} & \textbf{ResNet-LSTM} & \textbf{GraphSAGE-LSTM} \\
    \midrule
    \texttt{base} / \texttt{hidden\_g}           & $\{16,24,32,48\}$                     & $\{64,96,128,192,256\}$               \\
    \texttt{emb\_dim} / \texttt{n\_sage\_layers} & $\{64,128,192,256\}$                  & $\{2,3,4\}$                           \\
    \texttt{hidden\_t} (LSTM)                    & $\{64,128,192,256\}$                  & $\{64,96,128,192,256\}$               \\
    \texttt{n\_lstm\_layers}                     & $\{1,2\}$                             & $\{1,2\}$                             \\
    \texttt{k\_neighbors}                        & ---                                   & $\{4,8,12,16\}$                       \\
    \texttt{input\_bn}                           & ---                                   & $\{\text{True},\text{False}\}$        \\
    \texttt{concat\_agg}                         & ---                                   & $\{\text{True},\text{False}\}$        \\
    \texttt{l1\_reg}                             & $\{0,\,10^{-5},\,10^{-4},\,10^{-3}\}$& $\{0,\,10^{-5},\,10^{-4},\,10^{-3}\}$\\
    \texttt{dropout}                             & $\{0,0.1,0.2,0.3\}$                  & $\{0,0.1,0.2\}$                       \\
    \texttt{lr}                                  & $[3{\times}10^{-4},\;3{\times}10^{-3}]$ & $[5{\times}10^{-5},\;2{\times}10^{-3}]$ \\
    \texttt{weight\_decay}                       & $[10^{-6},\;10^{-3}]$                & $[10^{-6},\;10^{-3}]$                 \\
    \texttt{batch\_size}                         & $\{16,32,64\}$                        & $\{8,16,32\}$                         \\
    \bottomrule
  \end{tabular}
\end{table*}

The Optuna study identifies the best trial on validation $\rmseday$; the
winning configuration is then retrained from scratch with a longer epoch
budget and increased patience (ResNet-LSTM and GraphSAGE-LSTM: 30~epochs,
patience~8; FlatMLP: 35~epochs, patience~10), early-stopped on the same
validation split, and evaluated on the held-out test set. Because the
20-trial (50-trial for FlatMLP) Optuna budget is run independently
\emph{per seed} (Section~\ref{ssec:optuna}), the four-seed statistics
reported in Section~\ref{sec:results} reflect variability from both
independent hyperparameter search outcomes and training stochasticity
(weight initialisation, data shuffling) jointly, not training stochasticity
alone under a single fixed hyperparameter configuration.

\subsection{Training protocol}

All deep-learning models are trained with the Adam optimiser minimising the MSE
loss on normalised targets ($y_\text{norm} = (y - \bar{y}) / s_y$, with $\bar{y}$
and $s_y$ estimated from the training set). Training uses early stopping with
patience~6 during Optuna trials for ResNet-LSTM and GraphSAGE-LSTM (patience~7
for FlatMLP) and the higher final-retraining patience noted above. Automatic
Mixed Precision (AMP) is enabled when a GPU is available, except for SGLD
posterior sampling (Section~\ref{ssec:uq}), which is run entirely in fp32 so
that the injected Langevin noise is not distorted by mixed-precision loss
scaling. Each experiment is replicated with four seeds
$\{42, 1, 7, 13\}$ to report mean~$\pm$~std statistics, except the weighted
$k$-NN ablation, which uses two seeds $\{42, 1\}$ as noted in
Section~\ref{ssec:optuna}.

\subsection{Uncertainty quantification: epistemic uncertainty via SGLD posterior sampling}
\label{ssec:uq}

The best-performing architecture and hyperparameter configuration identified
in Sections~\ref{ssec:resnet}--\ref{ssec:mlp} is retrained from a fresh
initialisation under an SGLD posterior sampler, yielding a probabilistic
model that quantifies \emph{epistemic}
uncertainty \cite{kendall2017}: the reducible share of predictive uncertainty
due to limited or atypical training data and finite model capacity, as
opposed to the irreducible stochasticity of cloud formation itself
(\emph{aleatoric} uncertainty, which we do not model separately here; see
Section~\ref{sec:discussion}, Limitations). A point forecast alone cannot
distinguish a well-supported prediction from one made under conditions the
model has rarely seen during training; the latter case specifically signals
when a forecast should be treated cautiously (Section~\ref{sec:intro}).

\textbf{Epistemic uncertainty: SGLD posterior sampling.}
Standard SGD/Adam training returns a single point estimate $\hat\theta$ of
the network weights. To approximate the posterior distribution
$p(\theta \mid \mathcal{D}) \propto \exp(-\mathcal{L}(\theta))\,p(\theta)$
instead, we use Stochastic Gradient Langevin Dynamics (SGLD,
\cite{welling2011}) under an isotropic Gaussian prior
$p(\theta) = \mathcal{N}(0, \lambda^{-1}I)$ of fixed precision $\lambda$,
which injects calibrated Gaussian noise into each confined gradient step:
\begin{equation}
  \theta_{t+1} = \theta_t - \epsilon_t \bigl(\nabla_\theta \mathcal{L}(\theta_t)
  + \lambda\,\theta_t\bigr) + \sqrt{2\epsilon_t}\,\eta_t, \qquad
  \eta_t \sim \mathcal{N}(0, I),
  \label{eq:sgld}
\end{equation}
where $\epsilon_t$ is the (decaying) step size and the $\lambda\theta_t$ term
is the gradient of the negative log-prior, confining the chain around the
origin rather than allowing an unconfined random walk. Under suitable
conditions on $\epsilon_t$, the resulting Markov chain asymptotically samples
from $p(\theta \mid \mathcal{D})$ without the computational overhead of full
MCMC. We fix $\lambda = 100$ throughout, chosen independently of the
Optuna-tuned Adam weight decay used to obtain $\hat\theta$ in
Sections~\ref{ssec:resnet}--\ref{ssec:mlp}: the confined chain's stationary
variance scales as $1/\lambda$ and its relaxation timescale as
$1/(\epsilon_t\lambda)$ steps, so reusing Adam's much smaller weight decay
(${\sim}10^{-6}$) leaves the chain effectively unconfined and numerically
divergent over the sampling budget used here, which we verified empirically
before fixing $\lambda$ at this value. Training starts from a fresh random
initialisation (using the architecture and hyperparameters of the winning
Optuna trial, not its trained weights) and runs under the SGLD update rule
throughout; following a burn-in period of $T_\text{burn-in} = 500$ epochs
(discarded, while the chain moves from its random initialisation toward the
confined posterior regime), we collect $T = 10$ weight snapshots
$\{\theta_t\}_{t=1}^{T}$ at a fixed sampling interval of $\Delta = 100$
epochs, yielding an inexpensive approximate ensemble.

\textbf{Ensemble point forecast and epistemic variance.}
For a test input $x$, each posterior sample $\theta_t$ produces a point
prediction $\mu_t(x)$. The ensemble mean
$\bar\mu(x) = \frac{1}{T}\sum_{t=1}^T \mu_t(x)$ is the reported point
forecast, and the spread of the $T$ posterior predictions around it is the
epistemic variance:
\begin{equation}
  \sigma^2_\text{epi}(x) =
  \frac{1}{T}\sum_{t=1}^{T}\bigl(\mu_t(x)-\bar\mu(x)\bigr)^2.
  \label{eq:epistemic}
\end{equation}
$\sigma^2_\text{epi}(x)$ is large where the $T$ posterior samples disagree on
$\mu_t(x)$ --- i.e.\ where training data was scarce or atypical --- and
shrinks as more representative training data becomes available. We report
$\sigma_\text{epi}$ alongside the predictive interval
$\bar\mu(x) \pm z_{1-\alpha/2}\,\sigma_\text{epi}(x)$, evaluating empirical
coverage at nominal levels $1-\alpha \in \{80\%, 90\%, 95\%\}$. Because no
separate aleatoric term is estimated, this interval reflects only the
model's epistemic (weight-posterior) disagreement and should not be
interpreted as a full predictive interval capturing irreducible cloud
stochasticity as well (Section~\ref{sec:discussion}, Limitations).

\subsection{Toward full uncertainty disentanglement: a cooperative aleatoric
variance network}
\label{ssec:variance}

The SGLD pipeline above quantifies epistemic uncertainty alone. Training a
single network to output both a point forecast and its variance jointly is
known to produce imbalanced gradients and overconfident variance estimates
\citep{yi2025cooperative}; we are therefore extending the pipeline with the
cooperative, sequential strategy of \citet{yi2025cooperative}, which trains
a mean network, a variance network, and a Bayesian network in strict
sequence rather than jointly. The mean network of this strategy is already
the Optuna-tuned backbone of Sections~\ref{ssec:resnet}--\ref{ssec:mlp};
here we describe the second component, an aleatoric variance network,
currently being trained across the site-horizon-seed combinations for
which a mean-network checkpoint is available.

\begin{figure}[t]
  \centering
  \includegraphics[width=\columnwidth]{fig_uq_pipeline}
  \caption{Cooperative aleatoric/epistemic uncertainty pipeline. Step~1
           (mean network) is the already-trained, frozen backbone of
           Sections~\ref{ssec:resnet}--\ref{ssec:mlp}. Step~2 (variance
           network, this section) is implemented and currently being
           trained across the available backbones. Step~3 reuses the SGLD
           sampler of Section~\ref{ssec:uq}, but warm-started from the mean
           network's weights and re-weighted by the fixed $\sigma^2_\text{a}(x)$
           from Step~2 rather than the plain-MSE likelihood currently
           implemented there; this integration is not yet done
           (Section~\ref{sec:discussion}, Limitations).}
  \label{fig:uq-pipeline}
\end{figure}

\textbf{Aleatoric variance network.} With the mean network $\hat\theta$
from Sections~\ref{ssec:resnet}--\ref{ssec:mlp} frozen, we train a separate
network of the same architecture family whose final layer is replaced by a
two-unit head producing $(\alpha(x), \lambda(x))$, both mapped through a
softplus to ensure positivity. Following the Gamma-likelihood
reparameterisation of \citet{yi2025cooperative}, the network is trained to
predict the squared residual $r(x) = (\hat\mu(x) - y)^2$ of the frozen mean
network under a Gamma$(\text{shape}=\alpha(x), \text{rate}=\lambda(x))$
likelihood, minimising the negative log-likelihood
\begin{equation}
  \mathcal{L}_\text{var} = \mathbb{E}\bigl[-\alpha(x)\log\lambda(x)
  + \log\Gamma(\alpha(x)) - (\alpha(x)-1)\log r(x) + \lambda(x)\,r(x)\bigr],
  \label{eq:gamma_nll}
\end{equation}
which yields the aleatoric variance as the Gamma distribution's mean,
$\sigma^2_\text{a}(x) = \alpha(x)/\lambda(x)$. Training only on daytime
samples ($\text{GHI} \geq 20$~W/m²) keeps the variance network's target
consistent with the daytime evaluation convention used throughout this
paper. This component has been implemented and validated --- synthetic
recovery of known Gamma parameters, and cross-checked residual statistics
against the frozen backbone's own reported error --- and is currently being
trained across the available backbones; integrating its output into a
warm-started, variance-weighted SGLD sampler, completing the third
component of the cooperative strategy, is ongoing work
(Section~\ref{sec:conclusion}).

\subsection{Evaluation metrics}

Following standard practice in solar forecasting \cite{coimbra2013,zhang2015}, we report:

\textbf{Daytime RMSE} ($\rmseday$, W/m²): RMSE computed only over test samples
with observed $\text{GHI} \geq 20$~W/m², excluding night-time samples where all
models trivially predict zero.

\textbf{Daytime skill score} ($\skillday$): normalised improvement over persistence,
\begin{equation}
  \skillday = 1 - \frac{\rmseday(\text{model})}{\rmseday(\text{persistence})}
  \label{eq:skill}
\end{equation}
A score of 0 indicates parity with persistence; positive values indicate improvement.